\documentclass[a4paper,10pt,twocolumn,uplatex]{jsarticle}
\usepackage{style/nislab,style/resume}

%---------------------------------------------------------------------
% レジュメ種別・日付設定（要変更）
% \type{} 1:修士論文諮問会 2:卒業論文発表会 3:月例発表会 4:研究室合同発表会
%---------------------------------------------------------------------
\type{3}
\year{2026}
\month{6}
\date{8}

%---------------------------------------------------------------------
% ページ番号設定（要変更）
%---------------------------------------------------------------------
\setcounter{page}{9}

%---------------------------------------------------------------------
% 変更不要
%---------------------------------------------------------------------
\begin{document}

%---------------------------------------------------------------------
% タイトル作成部分（要変更）
% \maketitle{タイトル}{title}{名前}{name}
%---------------------------------------------------------------------
\maketitle{車両合流時におけるNR-V2Xユニキャスト利用時の通信信頼性向上の検討}
{Enhancing Communication Reliability of NR-V2X Unicast in Vehicle Merging}
{森 梓恩}
{Shion Mori}

%---------------------------------------------------------------------
\section{はじめに}
\label{はじめに}
自動運転の高度化に伴い高信頼・低遅延を満たす通信の実現が求められる．そこで車両同士が直接通信を行う車車間通信の研究が行われている．
車車間通信の規格として，LTE技術を応用したLTE-V2X PC5や5G技術を応用したNR-V2X PC5がある\cite{Tutorial}．
両規格はブロードキャストが可能であるが，NR-V2X PC5では1対1通信に特化した通信方式であるユニキャストが可能である．
ユニキャストは受信確認や再送制御が可能であるため，ブロードキャストより高い通信の信頼性が期待できる．
車車間通信を行うユースケースとして，高速道路の合流可能地点付近で，合流車線を走行する合流車両と本線車線を走行する本線車両が協調的に合流を行う合流時がある．以降，協調的に合流を行う複数車両を合流協調車両と定義する．
合流協調車両間の通信では，通常時の車両間と比較して高信頼・低遅延な通信要件が求められる．
そこで合流協調車両間のユニキャストが期待できる．
車車間通信において，通信を行うためにパケットを送信する時間と周波数を表す無線リソースを車両が自律的に選択する．
車車間通信の無線リソース割り当て方式としてSPS（Semi-Persistent Scheduling）方式がある．
SPS方式では，他車両が使用する無線リソースや受信電力を取得して，自車両の無線リソースの選択候補から除外することで，信頼性の高い無線リソースを選択できる．
しかし，SPS方式は各車両が互いの通信要件を満たすための協調的な無線リソース割り当てを行わないという課題がある．
そのため合流時のように異なる通信要件の車両が混在する環境の場合，各車両の通信要件を満たせない可能性がある．
本研究では，合流協調車両間でユニキャストを行う無線リソースを予約，共有する新たな無線リソース割り当て方式を検討する．
合流時の通信要件を考慮した協調的な無線リソース割り当てが可能になり合流協調車両間の通信信頼性向上を実現する．
%---------------------------------------------------------------------
\section{関連研究}
\subsection{無線リソース使用情報の共有}
Nikoorooらは，SPS方式の課題である車両同士のパケット衝突を軽減するために，自車両が検知した無線リソース使用情報を他車両と共有する無線リソース割り当て方式を提案した．
この手法により，SPS方式と比較して車両同士のパケット衝突が減少することで通信の信頼性が向上することを示した\cite{NR_RMAP}．
本研究では合流協調車両間の通信信頼性向上のために，本手法を応用して予約した無線リソース情報を他車両と共有する手法を検討する．

\subsection{緊急車両のための無線リソース割り当て方式}
Suranjanらは，緊急車両の通信の信頼性向上を目的とした無線リソース割り当て方式を提案した\cite{Emerge}．
提案方式では緊急車両が送信周期ごとに複数回送信を行う中で，送信を行わない時刻を設定する．
緊急車両はその時刻においてパケットの衝突が発生しているか検知できるため，緊急車両の継続的なパケット衝突を回避することができ，緊急車両の通信の信頼性が向上した．
しかし，合流時においては送信を行わない時刻が増加すると合流時の高信頼・低遅延な通信要件を満たさないという課題がある．


%---------------------------------------------------------------------
\section{提案方式}
\subsection{概要}
本研究では，合流車両の合流先車線上かつ前後に最も近い前方車両および後方車両の前後車両と，合流車両間が協調的に合流を行う環境を想定する．合流協調車両間のユニキャストにおける通信の信頼性向上を目的として，NR-V2X PC5 のSPS方式を拡張した新たな無線リソース割り当て方式を検討する．
合流車両の合流先車線上かつ後方に最も近い後方車両は合流協調車両間のユニキャストに用いる無線リソースを予約する．
後方車両は予約した無線リソースの予約情報を他車両に共有する．
これによりユニキャストを行わない車両が予約された無線リソースを避けて通信を行うため，ユニキャストを行う合流協調車両間の通信の信頼性向上が期待できる．

\subsection{合流前の動作手順}
合流車両の合流前において，後方車両は合流協調車両間でユニキャストを行うための無線リソースを予約，共有する．前後車両の決定手順および後方車両による無線リソースの予約と共有手順の2つの手順に従い実現する．

\subsubsection{前後車両の決定手順}
本線車両は，送信周期Rで自車両の位置情報と車両識別IDを含む車両情報をブロードキャストする．
合流車両は，本線車両の車両情報を取得する．
合流車両は，自車両の位置情報を用いて，現在の位置から合流可能地点までの距離が$L[m]$以内かを確認する．
合流車両が$L[m]$以内に位置するとき，合流車両は自車両の位置情報と本線車両の位置情報を用いて，前後車両をそれぞれ決定する．

\subsubsection{ユニキャスト用無線リソースの予約と共有手順}
合流車両は，前後車両を決定した後，各車両に対して合流時ユニキャスト対象車両であること示すユニキャスト対象通知メッセージを送信する．
さらに合流車両は，後方車両に対して合流協調車両間のユニキャストに用いる無線リソースの予約を要求する無線リソース予約要求メッセージを送信する．
後方車両は，ユニキャスト対象通知メッセージと無線リソース予約要求メッセージを受信した後，SPS方式により合流協調車両間のユニキャスト用無線リソースを選択する．
後方車両は，選択したユニキャスト用無線リソース情報を車両情報に追記して，送信周期Rで他車両にブロードキャストする．
後方車両からユニキャスト用無線リソース情報を受信した車両は，ユニキャスト用無線リソース情報を自車両の車両情報に追記して送信周期Rで他車両にブロードキャストする．

\subsection{合流時の動作手順}
合流車両は合流時，図\ref{fig:resource_selection}に示すように前方車両間および後方車両間のユニキャスト用無線リソースを用いて，それぞれ送信周期$R$に従いユニキャストを行う．

\begin{figure}[tb]
  \centering
  \includegraphics[width=\linewidth]{img/master_vehicleMerge.png}
  \caption{前方車両及び後方車両への無線リソース割り当て例}
  \label{fig:resource_selection}
\end{figure}





%---------------------
\section{評価}
\subsection{シミュレーション条件}
本研究ではネットワークシミュレータであるns-3（Net-work Simulator 3）を用いてシミュレーションを行う．
シミュレーションにおけるパラメータを\tabref{table:data_type}に示す．
高速道路片側3車線，合流車線の合流可能地点付近において，本線車両が$N$台存在し合流車両が1台が合流する環境で評価する．
車両間で車車間通信を行い，合流協調車両間の通信信頼性，およびその他本線車両の通信信頼性を評価する．
提案方式の有効性を評価するためにLTE-V2X PC5によるブロードキャスト，NR-V2X PC5によるブロードキャストおよびユニキャストと提案方式を比較する．
LTE-V2X PC5およびNR-V2X PC5によるブロードキャストでは，ユニキャストを行わず，ブロードキャストのみを用いて合流時の通信を行う条件で評価する．
\begin{table}
  \centering
  \caption{シミュレーションにおけるパラメータ値}
  \label{tab:delay}
  \begin{tabular}{lcr}
    \hline
      パラメータ & 設定\\
    \hline \hline
      車両台数N & [10，20，...500]台 \\
      合流可能地点からの距離L & 100m \\
      パケット送信周期$R$ & 100ms \\
      合流時パケット許容遅延時間 $PDB_{merge}$ & 10ms \\
      通常時パケット許容遅延時間 $PDB_{normal}$ & 100ms \\
      
    \hline
  \end{tabular}
  \label{table:data_type}
\end{table}



\subsection{評価指標}
通信の信頼性を評価するために，合流協調車両間の通信について，$PDB_{merge}$以内に受信成功した割合を示すパケット受信率$PDR_{merge}$を評価する．本線車両を走行する本線車両間の通信について，$PDB_{normal}$以内に受信したパケット受信率$PDR_{normal}$を評価する．また，パケットの持続的な衝突について評価するためにPIR（Peak-to-average Improvement Ratio），および平均AoI（Age of Information）を評価する．



\section{まとめ}
高信頼・低遅延を満たす通信の実現に向けて車車間通信におけるユニキャストが注目されている．
NR-V2X PC5ではユニキャストが可能であり，SPS方式により無線リソースを割り当てる．
しかし，SPS方式では各車両が互いの通信要件を満たすための協調的な無線リソース割り当てを行わない．
そのため合流時の合流協調車両間の高信頼・低遅延な通信要件を満たせない可能性がある．
本研究では合流協調車両間がユニキャストに用いる無線リソースを予約，共有する新たな無線リソース割り当て方式を検討した．
提案方式は従来方式と比較して，合流協調車両間の通信の信頼性向上に寄与する．


% Bibliography（参考文献）
%---------------------------------------------------------------------
% thebibliography を利用する場合は以下を使用
\footnotesize{
  \begin{thebibliography}{99}
    \bibitem{Tutorial}M. H. C. Garcia et al., "A Tutorial on 5G NR V2X Communications," in IEEE Communications Surveys Tutorials, vol. 23, no. 3, pp. 1972-2026, thirdquarter 2021
    \bibitem{NR_RMAP}S. Nikooroo, J. Estrada-Jimenez, A. Machalek, J. Härri, T. Engel and I. Turcanu, "Mitigating Collisions in Sidelink NR V2X: A Study on Cooperative Resource Allocation," 2024 22nd Mediterranean Communication and Computer Networking Conference (MedComNet), Nice, France, 2024, pp. 1-4
    \bibitem{Emerge}Suranjan Daw, Anwesha Kar, and Bheemarjuna Reddy Tamma. 2023. On Enhancing Semi-Persistent Scheduling in 5G NR V2X to Support Emergency Communication Services in Highly Congested Scenarios. In Proceedings of the 24th International Conference on Distributed Computing and Networking (ICDCN '23). Association for Computing Machinery, New York, NY, USA, 245–253.
  \end{thebibliography}
}

% BibTex を利用する場合は以下を使用（初めての人には難しいかも）
% \bibliographystyle{junsrt}
% \bibliography{myref}

%---------------------------------------------------------------------
\end{document}
%---------------------------------------------------------------------
